\itemhead{3.2 例题 · 例题1（教材p.81）}{XJ-B1-C03-EX-P081-Y0375}
\label{XJ-B1-C03-EX-P081-Y0375}
{\footnotesize 湘教版教材p.81 / PDF第88页 / 3.2 例题 / 例题1}
\field{原题}{已知函数 $f(x)=3x+b\ (x\in\mathbb R)$，判断该函数的单调性。}
\field{参考答案}{在 $\mathbb R$ 上单调递增，因为 $f(x_2)-f(x_1)=3(x_2-x_1)>0$。}
{\small \field{关联知识}{F05 单调性定义、证明与参数（教材帮p.97、98、100 / PDF106、107、109）}
\field{方法与易错点}{用定义比较任意两值，或依据已知图象性质；含参时使指定区间被相应单调区间覆盖。 易错：用有限采样代替证明；把分离半轴合并；漏跨零点情形。}
\field{答案依据}{例题原解位置：湘教版教材p.81及题库原图；独立编写可编辑参考解答。答案册此部分主要收练习、习题，不伪造例题答案册位置。}}
